% Full manuscript draft for structure and writing.
% This file is intentionally separate from generated tables/macros.
% Result values are pulled through generated macros where possible.

\documentclass[conference]{IEEEtran}

\usepackage{booktabs}
\usepackage{xspace}
\usepackage{amsmath}
\usepackage{graphicx}
\usepackage{cite}
\usepackage{url}

% Allow compilation either from the repository root or from the paper/
% directory. Generated numbers still come from the same provenance pipeline.
\newcommand{\PaperRoot}{}
\IfFileExists{paper/snippets/result_macros.tex}{%
  \renewcommand{\PaperRoot}{paper/}%
}{%
  \renewcommand{\PaperRoot}{}%
}
\newcommand{\PaperInput}[1]{\input{\PaperRoot#1}}
\newcommand{\PaperAsset}[1]{\PaperRoot#1}

\PaperInput{snippets/result_macros.tex}

\title{Native MXFP4 Persistent-State Dataflow for Gated DeltaNet Decode on FPGA}

\author{Anonymous Submission}

\begin{document}
\maketitle

\begin{abstract}
Linear-attention models reduce the asymptotic cost of long-context inference, but
single-token autoregressive decode remains constrained by the movement and update
of recurrent state. This paper studies one Gated DeltaNet decode layer from the
Qwen3-Next family and presents a native microscaled FP4 datapath for FPGA
implementation. The design stores persistent recurrent state on chip and uses
block-scaled MXFP4 operands, with E2M1 elements and shared E8M0 block scales, to
reduce memory traffic and arithmetic area while preserving the structure of the
delta-rule update. Unlike approaches that dequantize FP4 values into a wider
floating-point format before multiplication, the inner multiply is implemented as
native integer and LUT logic for sign, exponent, mantissa, normalization, and
fixed-point accumulation. On an Alveo U55C target, the default configuration
reaches \OursMxfpFourBThirtyTwoLatencyUs per token in RTL cosimulation and
\OursMxfpFourBThirtyTwoPowerW post-implementation power, yielding
\HeadlineSpeedupVsHOneHundred speedup over the cited H100 baseline and
\HeadlineSpeedupVsUsc speedup over the cited USC BF16 FPGA baseline. The current
artifact includes deterministic synthetic accuracy checks; realistic Qwen
activation capture remains \QwenCaptureStatus and is therefore not claimed as a
completed PPL result in this draft.
\end{abstract}

\section{Introduction}

Autoregressive decoding in large language models introduces measurable latency at
inference time because each generated token requires sequential access to model
state from previous tokens. It is often described as a matrix-multiply problem,
but the decode step of modern sequence models also has a persistent-state
problem. In linear-attention layers, each new token updates a compact recurrent
state rather than appending to a full attention cache. This improves the
asymptotic memory profile, yet it changes the hardware bottleneck: the
accelerator must repeatedly read, modify, and preserve a structured state object
with low latency.

Gated DeltaNet (GDN) layers are an especially useful target for this regime. A
single decode step combines a key-value state lookup, a delta-rule update, and an
output gate. The resulting computation is small enough that launch overheads and
memory movement matter, but regular enough that an FPGA can exploit persistent
storage and customized arithmetic. Prior FPGA work on persistent-state dataflow
for linear attention demonstrated the value of keeping state close to the
compute fabric~\cite{usc_linear_attention_fpga}, but used BF16 arithmetic for
the core datapath. BF16 is robust, but it is not obviously the best point in the
design space when the operation is decode-only, batch size is fixed, and the
state is accessed every token.

This paper asks whether microscaled FP4 can close the gap between aggressive
quantization and model-quality preservation for this setting. Flat INT4 is
attractive for area, but it has limited dynamic range and is sensitive to
outliers. MXFP4 instead pairs compact E2M1 elements with an E8M0 scale shared by
a block of elements. This gives the datapath a small floating-point-like dynamic
range without paying the cost of BF16 multipliers. The key hardware question is
whether those block scales can be handled natively, inside the accumulation
path, without expanding the inner loop back to FP32 or BF16. The format follows
the microscaling direction standardized by OCP and reflected in the Microsoft
\texttt{microxcaling} reference implementation~\cite{ocp_mx,microxcaling}.

The contributions of this research are fivefold. First, we provide a native
MXFP4 persistent-state datapath for a batch-one Gated DeltaNet decode layer on
an FPGA. Second, we implement the FP4 E2M1 element multiply in LUT logic with
shared block-exponent handling. Third, we specialize the token dataflow so the
recurrent state remains resident during decode. Fourth, we evaluate latency,
timing, resource utilization, power, and synthetic numerical fidelity through
HLS C/RTL cosimulation and Vivado post-implementation reports. Fifth, we report
recurrent-state stress and state-storage/traffic accounting while reserving
full-model quality claims for future realistic activation capture.

\section{Background and Motivation}

\subsection{Gated DeltaNet Decode}

For a single token, GDN maintains a recurrent state matrix per head. Given query
$q_t$, key $k_t$, value $v_t$, learning-rate parameter $\beta_t$, output gate
$g_t$, and prior state $S_{t-1}$, the reference update is

\begin{equation}
S_t = S_{t-1} - \beta_t (S_{t-1} k_t - v_t) k_t^\top.
\end{equation}

The output is produced by querying the updated state and applying the gate. The
important hardware property is that the state is both read and updated every
token. If this state is streamed through external memory on each invocation, the
kernel becomes dominated by interface traffic rather than arithmetic. Our final
default kernel therefore keeps the recurrent state resident in partitioned
on-chip memory and exposes only token inputs and packed token outputs at the
top-level interface.

\subsection{Algorithmic Mapping}

The hardware mapping follows the algebra of the delta rule but changes its
storage and scheduling. Conceptually, each token performs four operations.
First, the current key is used to read the prior state and form the predicted
value. Second, that prediction is compared with the incoming value to form a
residual. Third, the residual and key form a rank-one correction, scaled by
$\beta_t$. Fourth, the corrected state is queried and gated to produce the token
output. The implementation keeps these operations in a fixed token schedule so
that state read, correction, and writeback are visible to HLS as a steady
pipeline.

This mapping is useful because it separates two forms of parallelism. The
key-direction parallelism controls how many columns of the state participate in
one update step, while the value-direction parallelism controls how many output
lanes are updated together. Increasing either factor without matching memory
banking can create conflicts, so the design treats parallelism and state layout
as a single choice rather than independent parameters.

\subsection{MXFP4 Quantization Method}

The proposed design utilizes MXFP4 as the primary arithmetic format. MXFP4 is a
representation where each value is represented as a 4-bit E2M1 element (one sign
bit, two exponent bits and one mantissa bit), along with a common E8M0 block
exponent. Unlike formats that assign an independent exponent to each value,
MXFP4 provides a shared block scale for all values within a block. Assigning one
shared block scale to all values within a block provides a reduction in memory
usage while providing greater dynamic range than would be possible with a flat
INT4 format.

Each value in a block is approximated as the product of the shared block scale,
and a quantized E2M1 element. The shared exponent captures the scale of the
entire block, while the E2M1 element captures the signed mantissa and exponent
information for each individual element.

The default hardware datapath operates with an MXFP4 block size of 32. The
software reference uses the same block grouping and rounding behavior as the HLS
implementation so that synthetic test vectors can be compared bit-exactly across
the Python model, C-simulation, and RTL cosimulation.

The FP4 element multiply is implemented as a native E2M1 operation, rather than
by converting operands into BF16 or FP32. The multiply path consists of: sign XOR
on the sign bits, exponent addition with bias adjustment, mantissa
multiplication, normalization; the normalized result of the element multiply is
then aligned based upon the shared block exponents, before being accumulated.

MXFP4 quantizes elements differently than flat INT4. Flat INT4 uses extremely
compact values, but it lacks local dynamic-range flexibility. MXFP4 maintains
the same 4-bit element width while using block exponents to adapt to local
tensor magnitude.

\subsection{Why FPGA}

FPGAs are not a replacement for GPUs in full-model serving, but they are useful
for studying dataflow choices that GPUs cannot express directly. A custom FPGA
kernel can keep persistent state physically close to the arithmetic, avoid DSPs
for unsupported low-precision formats, and specialize interfaces to the exact
decode path. This makes FPGA a good testbed for the claim that MXFP4 can provide
a more useful accuracy-area tradeoff than flat INT4 for linear-attention decode.

\section{Related Work}

\subsection{Linear-Attention Accelerators}

Prior accelerators for transformer inference have focused primarily on dense
matrix multiplication, KV-cache movement, and prefill throughput. Linear
attention changes the shape of the problem by replacing the token-length cache
with a recurrent state. The closest baseline for this work is the USC
persistent-state dataflow accelerator for memory-bound linear attention decode,
which demonstrates that keeping state close to the FPGA fabric can outperform a
general GPU path for batch-one decode. This paper builds on that insight but
changes the arithmetic question: rather than asking whether BF16 stateful
linear-attention decode is efficient, it asks whether an MXFP4-native datapath
can reduce the arithmetic and memory footprint without losing the structure of
the update.

\subsection{Low-Precision Floating Point}

Low-precision floating point is attractive because it preserves some dynamic
range while reducing storage and multiplier cost. BF16 and FP16 are now common
in accelerator datapaths, but four-bit formats are less mature, especially on
FPGA. Flat INT4 is simple, but its single scale is often brittle for activations
and recurrent state. Microscaling formats address this by sharing an exponent
over a small block of elements. The result is a middle point between fixed-point
quantization and full per-element floating point: the element path is tiny, but
the block scale preserves local magnitude.

\subsection{FPGA Support for Non-Native Formats}

Modern FPGA DSP blocks are optimized for integer and conventional floating-point
building blocks, not for E2M1 FP4. A design that maps FP4 arithmetic through DSP
macros risks losing the very area advantage that motivates the format. This is
why the multiplier primitive in this work is explicitly LUT-based. The broader
design uses HLS for integration and dataflow exploration, but the arithmetic
primitive and block-scale alignment are written to expose the native FP4
structure rather than relying on a dequantize-then-multiply implementation.

\section{Design}

\subsection{Kernel Organization}

The hardware kernel follows five conceptual phases: input preparation, state
read, delta-rule update, state writeback, and output projection. In the
implemented default path, the state read and writeback operate on a resident
state tensor rather than a full external-memory state stream. This choice is the
main architectural change needed to reach the Phase 6 hard-floor latency target.
The output path is packed by the value parallelism factor, so a group of output
lanes is emitted as one wider word rather than as scalar stores.

The default compile-time configuration uses the MXFP4 block-size and parallelism
choices summarized in Table~\ref{tab:main}. The design is parameterized over
key-direction parallelism, value-direction parallelism, and MX block size, which
allows the sweep in Section~\ref{sec:results}.

The top-level design is intentionally narrow in scope. It accelerates one layer,
one token, and batch size one. This restriction is not a weakness for the target
question; it isolates the decode bottleneck and makes the effect of persistent
state visible. A full serving system would wrap this kernel with projection
layers, scheduling, memory management, and model-parallel execution, but those
systems concerns would obscure the arithmetic and dataflow contribution.

\begin{figure}[t]
  \centering
  \includegraphics[width=\linewidth]{\PaperAsset{figures/gdn_dataflow.pdf}}
  \caption{Persistent-state GDN decode dataflow. Token inputs are prepared as
  MX blocks, the recurrent state remains resident in partitioned on-chip memory,
  and the delta update writes back in place before packed output emission.}
  \label{fig:dataflow}
\end{figure}

Figure~\ref{fig:dataflow} summarizes the resident-state token path and its
packed output interface.

\subsection{Token Schedule}

At the token level, the accelerator follows a deterministic schedule. The input
phase quantizes or accepts already-packed token vectors and associates them with
block scales. The prediction phase reads the relevant state banks and produces
the state-vector product used by the delta rule. The update phase computes the
residual and rank-one correction. The writeback phase commits the corrected
state to the resident banks. Finally, the output phase packs multiple value
lanes into the top-level output word. This schedule is intentionally simple:
there is no dynamic control path, no data-dependent memory traversal, and no
token reordering.

The schedule also explains why later optimization work should focus on token
ingress and phase balance rather than simply increasing arithmetic parallelism.

\subsection{Native E2M1 Multiply}

The E2M1 multiply primitive decomposes each operand into sign, exponent, and
mantissa fields. The sign path is a bitwise XOR. The exponent path adds the small
exponents and applies the format bias. The mantissa path handles the implicit
leading bit for normal values and the mantissa-only interpretation of subnormal
values. The product is normalized and rounded into a fixed Q4.3 partial product.
This output format is narrow enough to keep the primitive LUT-based, but wide
enough to feed the block accumulator without immediately losing the local
magnitude information from the MX scale.

The key point is that this multiply is not implemented by widening the operands
into FP32 or BF16. That distinction matters for the area story: the FP4 path is
native in the arithmetic inner loop, while the accumulation path uses integer and
fixed-point operations with scale-aware alignment.

\begin{figure}[t]
  \centering
  \includegraphics[width=\linewidth]{\PaperAsset{figures/mxfp4_mac_datapath.pdf}}
  \caption{Native MXFP4 arithmetic path. E2M1 element multiplication is kept in
  LUT logic, while shared E8M0 block exponents are carried into the fixed-point
  alignment and accumulation stage.}
  \label{fig:mxfp4mac}
\end{figure}

Figure~\ref{fig:mxfp4mac} marks the boundary between ordinary quantization and
the paper's arithmetic contribution: the element multiply remains narrow, and
hardware is spent where the block exponent affects accumulation.

\subsection{Block-Exponent-Aware Accumulation}

MX formats complicate accumulation because each group of elements has a shared
scale. The block-exponent alignment module receives a stream of partial products
and their corresponding operand block exponents. It computes an effective block
scale, aligns partial values to a common exponent, and accumulates into a
saturating fixed-point register. At a block boundary, it emits the accumulated
value with the dominant exponent for downstream conversion.

This approach preserves the hardware benefit of FP4 elements while respecting
the numerical role of the shared scale. It is also the part of the design most
closely tied to linear-attention state updates, because the state, key, value,
and gate tensors can have different local scale structure.

\subsection{Persistent State and Packed Output}

An earlier version of the implementation streamed the full recurrent state
through top-level memory ports on every token. That version was useful for
debugging, but it missed the latency hard floor. The current design replaces
that stream with a static resident state array bound to on-chip RAM and
partitioned cyclically across both row and column dimensions. State readback is
retained as an optional debug path rather than the default data path.

The output interface is similarly specialized. The default kernel stores
multiple value lanes into one output word, reducing the number of scalar output
transactions. This is a small interface change, but it has a visible effect on
the report: the packed-output design reduces the output phase to a short pass
relative to the predict and update phases.

\subsection{Design Invariants}

Several invariants are maintained across the implementation. First, the state
layout keeps E2M1 element storage separate from E8M0 block scales, which allows
the state memory to be partitioned independently from the scale memory. Second,
the inner multiply never materializes FP32 operands. Third, block exponents are
handled before accumulation boundaries are crossed, so downstream modules see a
fixed-point value with an associated scale rather than a bag of unrelated FP4
codes. Fourth, test vectors are generated deterministically and reused across
software, C-simulation, and cosimulation.

\subsection{Memory Layout and Banking}

The recurrent state has a different access pattern from a conventional dense
matrix operand. Each token reads the state for prediction, forms the delta
correction, and then writes the updated state back before the next token. A
layout optimized only for burst transfer is therefore the wrong objective. The
implementation stores the state in banks that expose parallel row and column
access to the update loop. E2M1 element storage and E8M0 scale storage are kept
separate, so element bandwidth and scale bandwidth can be balanced
independently. This matters because a single scale is reused across a block of
elements; fetching the scale every element would erase part of the MXFP4 storage
benefit.

The state banks are partitioned cyclically to match the key and value
parallelism factors. The goal is not simply to maximize memory ports, but to
avoid banking conflicts in the inner update. When a block of keys is processed,
the corresponding state columns and scale entries are available to the MAC
fabric without serializing the block. This is the practical reason that block
size is part of the architectural sweep rather than only a quantization
hyperparameter.

\subsection{Rounding, Saturation, and Debug Modes}

The software and hardware models both treat MXFP4 as an integer datapath with
explicit rounding and saturation points. Element products are rounded into the
Q4.3 partial-product format, and block accumulation saturates into the INT24
range. These boundaries make the HLS behavior auditable: if a C-simulation or
RTL-cosimulation mismatch appears, the debug path can compare the product,
aligned partial, block accumulator, and final packed output separately.

Several debug features are deliberately outside the headline datapath. State
readback can be enabled to inspect the resident state, and per-phase testbenches
can isolate preparation, update, writeback, and output behavior. These features
are important for development but are not counted as the optimized deployment
path. The paper reports the packed-output persistent-state configuration because
that is the design point matching the intended decode dataflow.

\section{Methodology}

\subsection{Experimental Questions}

The evaluation is organized around four questions. First, does the persistent
MXFP4 kernel meet the latency hard floor for single-token decode? Second, does
the LUT-based FP4 datapath keep FPGA resource use low enough to leave room for
system integration? Third, do block size and parallelism expose a useful Pareto
curve rather than a single brittle point design? Fourth, how much numerical error
appears under deterministic synthetic vectors?

\subsection{Golden Model and Quantization}

The software reference begins with a direct FP32 implementation of the GDN decode
equation. The MXFP4 model then applies block quantization to activations,
weights, and recurrent state, using deterministic synthetic vectors for the
current artifact. State MXFP8 is included as a fallback comparison. All random
vectors are seeded, and generated test vectors are content-addressed.

The current artifact does not include real Qwen3-Next activation captures; that
status is reported as \QwenCaptureStatus. The accuracy results therefore remain
synthetic rather than full-model PPL evidence.

\subsection{HLS and Vivado Flow}

The hardware is implemented in Vitis HLS and targeted to the Xilinx Alveo U55C
device. C-simulation checks bit-exact output parity against the local HLS
testbench. C/RTL cosimulation is used for authoritative latency. Vivado
post-implementation reports provide LUT, BRAM, DSP, timing, and power
measurements. The paper-number pipeline extracts all values from report files
into \texttt{paper/numbers.json}, then renders LaTeX tables, macros, and figures.

The default configuration is treated as the paper headline configuration. Sweep
points are still implemented and reported, but they are used to understand the
design space rather than to cherry-pick the fastest isolated number. This matters
because FPGA timing closure and routing pressure can make a high-parallelism
configuration less useful than it appears at the C-synthesis stage.

\subsection{Verification Hierarchy}

The verification flow is layered so that failures can be localized before they
reach the full kernel. Format-level tests check E2M1 encoding, block scaling,
subnormal handling, and adversarial all-maximum or all-minimum blocks. The
software GDN tests then compare the quantized model with the FP32 reference under
deterministic synthetic inputs. HLS C-simulation checks the same token paths at
the C++ level, and RTL cosimulation is used for the latency number reported in
the paper. This hierarchy prevents a convenient but dangerous failure mode: a
full-kernel test can pass on a benign vector even while a subnormal, saturation,
or block-boundary corner case is broken.

The current artifact treats the Python MXFP4 model as the bit-exact behavioral
reference for HLS. Once the RTL is generated, cosimulation becomes the
authoritative latency source. Vivado implementation is used only after
cosimulation agrees with the golden model, because a post-implementation timing
or power number has little value if the RTL no longer represents the intended
arithmetic.

\begin{table}[t]
  \centering
  \caption{Kernel phases and verification anchors. The table is qualitative;
  measured values remain in generated result tables.}
  \label{tab:phasecheck}
  \resizebox{\linewidth}{!}{%
  \begin{tabular}{lll}
    \toprule
    Phase & Main responsibility & Verification anchor \\
    \midrule
    Prepare & Pack token vectors and attach MX scales & MX primitive tests \\
    State read & Fetch resident E2M1 state and E8M0 scales & Per-phase C-sim \\
    Update & Form residual and scaled rank-one correction & Python/HLS parity \\
    Writeback & Commit corrected state in place & State trace comparison \\
    Output & Gate and pack value lanes & Full-kernel cosim \\
    \bottomrule
  \end{tabular}}
\end{table}

Table~\ref{tab:phasecheck} clarifies what each kernel phase proves. This matters
for reviewability: a single end-to-end pass/fail result is less useful than a
phase-localized chain of evidence, especially when the arithmetic format is
non-standard and the implementation intentionally avoids DSP-backed
floating-point operators.

\subsection{Baselines}

We compare against two cited baselines from the USC persistent-state linear
attention accelerator paper: an NVIDIA H100 PCIe BF16 row and a USC BF16 FPGA
row. These are external citation values, not measurements re-run in this
repository. Our default row is measured by HLS RTL cosimulation and Vivado
post-implementation for the packed-output persistent-state MXFP4 design.

\section{Results}
\label{sec:results}

\subsection{Main Comparison}

Table~\ref{tab:main} shows the headline comparison. The default MXFP4 design
achieves \OursMxfpFourBThirtyTwoLatencyUs per token in RTL cosimulation, with
\OursMxfpFourBThirtyTwoLutPct LUT utilization, \OursMxfpFourBThirtyTwoBramPct
BRAM utilization, and \OursMxfpFourBThirtyTwoPowerW post-implementation power.
This corresponds to \HeadlineSpeedupVsHOneHundred speedup over the cited H100
baseline and \HeadlineSpeedupVsUsc speedup over the cited USC BF16 FPGA row.

\begin{table}[t]
  \centering
  \caption{Main latency, area, and power comparison. Baseline rows are external
  citation values; our row is extracted from RTL cosimulation and Vivado
  post-implementation reports.}
  \label{tab:main}
  \resizebox{\linewidth}{!}{\PaperInput{tables/main_results.tex}}
\end{table}

The latency entries in Table~\ref{tab:main} are the most direct evidence for
the accelerator claim. The default design does not merely reduce arithmetic
precision; it changes where state lives and how outputs are emitted. Without
those dataflow changes, the FP4 datapath alone would not be sufficient to reach
the measured latency point.

\subsection{Timing, Area, and Power}

The HLS C-synthesis report estimates \OursMxfpFourBThirtyTwoCsynthFmaxMhz for
the default kernel. Post-implementation timing closes with
\OursMxfpFourBThirtyTwoImplWnsNs worst negative slack margin. DSP utilization is
\OursMxfpFourBThirtyTwoDspPct, even though the FP4 primitive itself is intended
to be LUT-based; remaining DSP use comes from wider fixed-point arithmetic
introduced by the HLS implementation. This is an important engineering detail:
the current artifact validates the dataflow and packed-output latency path, but
there remains room to further constrain HLS arithmetic mapping in a later
revision.

\begin{table}[t]
  \centering
  \caption{Area breakdown for the default MXFP4 configuration.}
  \label{tab:area}
  \PaperInput{tables/area_breakdown.tex}
\end{table}

Table~\ref{tab:area} shows that the default point is not close to the U55C
capacity limits. This is important for two reasons. First, it suggests that
future integration with adjacent projection or scheduling logic has resource
headroom. Second, it indicates that the remaining performance gap to the
aspirational target is more likely to come from interface and dataflow
refinement than from exhausting fabric resources.

\subsection{Parallelism and Block-Size Sweep}

The Phase 6 sweep evaluates several combinations of key parallelism, value
parallelism, and block size. The default row is the only current row intended for
headline latency. Non-default rows are retained as sweep evidence and are useful
for explaining area and routing sensitivity. In particular, the max-parallel
points consume substantially more LUT and DSP resources and should be treated as
Pareto data rather than as the default design point.

The sweep rows in Table~\ref{tab:main} also explain why the default
configuration is a reasonable paper point even when larger parallelism settings
exist. On an FPGA, parallelism has to pay for routing, memory banking, and
timing closure. The best design point is therefore not the one with the largest
number of lanes, but the one that keeps the full token path balanced.

\subsection{Quantization Accuracy}

The deterministic synthetic quantization tests show that MXFP4 preserves output
direction more closely than the INT4 fallback under the current synthetic
workload. For block size 32, synthetic output cosine is
\SyntheticMxfpFourBThirtyTwoOutputCosine and output relative L2 error is
\SyntheticMxfpFourBThirtyTwoOutputRelLTwo. The INT4 fallback output cosine is
\SyntheticIntFourOutputCosine. The block-size ablation in
Table~\ref{tab:block} compares block sizes under the same synthetic harness.

\begin{table}[t]
  \centering
  \caption{Synthetic block-size ablation. These values do not replace
  full-model quality measurements.}
  \label{tab:block}
  \PaperInput{tables/ablation_block_size.tex}
\end{table}

The per-tensor ablation in Table~\ref{tab:quant} gives a first-order view of
which tensors are most sensitive in the synthetic setting and should guide
later calibration work.

\begin{table}[t]
  \centering
  \caption{Synthetic per-tensor quantization ablation.}
  \label{tab:quant}
  \PaperInput{tables/ablation_quantization.tex}
\end{table}

\subsection{State Precision and Traffic}

The newer recurrent-state stress results clarify the quality boundary of the
current synthetic evidence. MXFP4 remains the native four-bit arithmetic path,
while MXFP8 recurrent state remains the safer fallback if accumulated
recurrent-state drift becomes the limiting accuracy factor.

\begin{table}[t]
  \centering
  \caption{Synthetic recurrent-state stress evidence.}
  \label{tab:state-stress}
  \resizebox{\linewidth}{!}{\PaperInput{tables/recurrent_state_stress.tex}}
\end{table}

The storage and traffic accounting in Table~\ref{tab:state-traffic} supports
the persistent-state dataflow claim. The compact MXFP4 state footprint allows
the recurrent state to remain resident on chip, eliminating steady-state
off-chip recurrent-state traffic in the evaluated decode path.

\begin{table}[t]
  \centering
  \caption{Recurrent-state storage and off-chip traffic accounting for MXFP4
  $B=32$.}
  \label{tab:state-traffic}
  \resizebox{\linewidth}{!}{\PaperInput{tables/state_storage_traffic.tex}}
\end{table}

\subsection{Roofline Positioning}

Table~\ref{tab:roofline} summarizes the roofline-style interpretation for the
U55C target. The purpose is not to claim a universal peak throughput number; the
kernel is intentionally narrow and stateful. Instead, the summary identifies
whether the implementation is limited mainly by arithmetic throughput, memory
traffic, or the serialized portions of the token pipeline.

\begin{table}[t]
  \centering
  \caption{Roofline-style interpretation of the measured default design.}
  \label{tab:roofline}
  \begin{tabular}{lp{0.56\linewidth}}
  \toprule
  Aspect & Interpretation \\
  \midrule
  State traffic & Recurrent state remains on chip \\
  Arithmetic path & Native MXFP4 LUT multiply \\
  Current limiter & Token ingress and HLS scheduling \\
  Next target & Packed or streamed token path \\
  \bottomrule
  \end{tabular}
\end{table}

For this artifact, the roofline view reinforces the dataflow conclusion. The
state-resident kernel moves the design away from a pure external-memory limit,
but the remaining scalar token ingress and HLS-generated fixed-point operations
leave room before the design reaches the idealized compute ceiling.

\section{Reproducibility and Artifact}

The artifact is designed so that paper numbers can be regenerated rather than
copied. HLS, Vivado, cosimulation, benchmark, and figure-generation outputs are
written to fixed repository paths. The benchmark stage extracts measurements
into \texttt{paper/numbers.json}; \texttt{paper/provenance.json} records the
source file, source line where available, extractor, timestamp, and git SHA for
each value. The LaTeX snippets and tables used in this draft are generated from
those files.

\begin{figure}[t]
  \centering
  \includegraphics[width=\linewidth]{\PaperAsset{figures/paper_pipeline.pdf}}
  \caption{Paper-number provenance pipeline. Reports feed extractors, extractors
  update canonical JSON files, and generated LaTeX assets are consumed by the
  manuscript and paper pack.}
  \label{fig:pipeline}
\end{figure}

This discipline affects the writing style. Numerical claims in the manuscript
are expressed through generated macros such as
\texttt{\textbackslash OursMxfpFourBThirtyTwoLatencyUs} rather than hand-written
values. When the implementation changes and the benchmark pipeline is rerun, the
abstract, tables, and results section update together. For a hardware paper,
this reduces the risk that a late timing, power, or latency change leaves stale
numbers in the narrative.

\section{Threats to Validity}

The primary internal threat is HLS mapping. Although the FP4 multiply primitive
is designed as LUT logic, HLS may still infer DSPs for surrounding fixed-point
operations unless constrained carefully. The current resource report should
therefore be interpreted as the resource usage of this implementation, not as a
lower bound for a hand-written RTL design.

The primary external threat is model representativeness. Synthetic vectors are
useful for bit-exact testing and stress cases, but they cannot replace real
activation distributions.

The baseline comparison is also limited. H100 and USC BF16 rows are cited values
from a prior paper, while our row is measured in this repository's flow. The
comparison is still informative because all rows target decode latency, but a
perfect apples-to-apples comparison would require rerunning all baselines under a
single harness and software stack.

\section{Discussion}

The main lesson is that native low-precision arithmetic only helps if the
surrounding dataflow lets it matter. The first implementation streamed the full
state through the top-level interface and therefore spent most of its time moving
state rather than using the FP4 datapath. The persistent-state rewrite changed
the bottleneck: predict, update, and token I/O are now visible at similar scale,
and the packed output interface removes a full scalar output pass from the
critical path.

The second lesson is that MXFP4 is a better hardware target than flat INT4 for
this kind of recurrent decode, at least structurally. The shared exponent gives
the state update more local dynamic range without requiring BF16 multipliers.
The price is more complicated accumulation, since block exponents must be
aligned and carried through the update. That price is acceptable in an FPGA
setting because the alignment logic can be made explicit and amortized across
the block.

The third lesson is about reproducibility. Hardware papers often suffer from
late-stage number drift, where a table, figure, and abstract silently describe
different builds. This artifact avoids that failure mode by generating the
reported values, tables, and macros from one JSON file with source-line
provenance. The paper text should preserve that discipline by using the generated
macros rather than typing measured values by hand.

\section{Limitations}

This is a simulation-only study. We do not report physical board measurements,
end-to-end model throughput, or full Qwen3-Next serving results. The hardware
kernel covers one GDN decode layer and batch-one single-token decode.

The most important missing item is realistic Qwen activation capture. Because
the status is \QwenCaptureStatus, current accuracy evidence is synthetic-only and
should not be framed as Qwen PPL preservation.

The current HLS implementation also leaves optimization headroom. The default
design clears the latency and speedup hard floors, but it does not yet hit the
aspirational latency target. Token-input ingress remains scalar, and HLS maps
some wider arithmetic into DSP resources. A hand-tuned RTL version of the
accumulator and a packed input interface are natural next steps.

\subsection{Future Work}

Future work should complete the realistic-data path, run full perplexity
evaluation on a standard benchmark, and route those measurements through the
same provenance pipeline used for the current hardware results. On the
implementation side, the next revisions should pack token inputs, further
constrain non-FP4 arithmetic mapping, and evaluate the design on physical FPGA
hardware to compare Vivado-estimated power with board-level measurements.

\section{Conclusion}

This paper draft presents a native MXFP4 persistent-state datapath for Gated
DeltaNet decode on FPGA. The implemented design combines LUT-based E2M1
multiplication, block-exponent-aware accumulation, resident recurrent state, and
a packed output interface. On the current U55C flow, the default configuration
achieves \OursMxfpFourBThirtyTwoLatencyUs per token in RTL cosimulation and
\HeadlineSpeedupVsHOneHundred speedup over the cited H100 baseline, while using
\OursMxfpFourBThirtyTwoPowerW post-implementation power. The result supports the
central architecture argument: MXFP4 can make sense for linear-attention decode
when the state dataflow is designed around persistent storage rather than
external-memory streaming.

\section*{Acknowledgments}

Omitted for anonymous review.

\begin{thebibliography}{9}

\bibitem{usc_linear_attention_fpga}
Gupta et al. A Persistent-State Dataflow Accelerator for Memory-Bound Linear
Attention Decode on FPGA. arXiv:2603.05931, 2026.

\bibitem{ocp_mx}
Open Compute Project. Microscaling Formats Specification, version 1.0.

\bibitem{microxcaling}
Microsoft. microxcaling reference implementation.

\bibitem{qwen_next}
Qwen Team. Qwen3-Next model family.

\end{thebibliography}

\end{document}
